% Proposed caption only. This file is not included in the manuscript.
\caption{Probabilistic predictions and empirical finite-shot VQE executions under a matched configuration. All results use the ideal-Born benchmark, the original grid and cyclic candidate rule, success defined by $E_{\mathrm{ideal}}(z)-E_{\mathrm{ref}}\leq 0.02$, and iteration bound $K=50$.
(a) Success probability comparison between Gaussian-DTMC propagation and cumulative first-success frequencies from 400 independently seeded finite-shot executions. Hollow diamonds denote independently recorded Storm bounded-reachability checkpoints. Shading denotes pointwise $95\%$ Wilson intervals. At $K=50$, the prediction is $92.80\%$, compared with $91.50\%$ ($366/400$; interval $[88.36\%,93.85\%]$).
(b) Empirical-minus-Gaussian discrepancies in percentage points; the shaded intervals translate the same Wilson limits relative to the fixed model, with an endpoint discrepancy of $-1.30$ percentage points. These intervals quantify sampling uncertainty, not Gaussian approximation error.
(c) Unconditional first-success masses aggregated from exact iteration-level model outputs and per-run records; empirical bar labels are counts out of 400. The separate no-success-by-50 category contains 34 budget-terminated executions, and each complete outcome partition sums to $100\%$.
(d) Model and empirical conditional mean iterations (upper) and physical shots (lower), restricted to success within 50 iterations. Squares denote model means and circles denote empirical means. Error bars show the original $95\%$ percentile-bootstrap intervals over 366 successful paths, with 5,000 resamples. Physical shots equal $9{,}216$ times iterations, so the lower plot is a rescaled view of the same resource comparison. The marginal intervals exclude the modeled means; the previously adopted post-audit comparison gives raw $p=0.024$ and Holm-adjusted $p=0.094$ across P1@30, P1@40, P1@50, and P2. Conditional model means are obtained by Gaussian-kernel propagation, while Storm verifies reachability checkpoints. The observations are statistically compatible at the evaluated resolution under this analysis; equivalence is not established.}
\label{fig:probabilistic-results-expanded}
